% =====================================================================
%  Anatomía de un LLM — Fase III — Hito 3
%  Examen de hito: el Transformer (módulos 12 a 15)
%  Compilar desde esta carpeta: latexmk -r ../latexmkrc H3-examen.tex
% =====================================================================
\documentclass[11pt,a4paper,twoside,openany]{book}
\usepackage{../estilo-apuntes}
\usepackage{float}

\newcommand{\ApuntesCurso}{Anatomía de un LLM}
\newcommand{\ApuntesModulo}{Hito 3}

\newcommand{\T}{^{\top}}
\newcommand{\dff}{d_{\mathrm{ff}}}
\newcommand{\RMSNorm}{\operatorname{RMSNorm}}
\newcommand{\MHA}{\operatorname{MHA}}
\newcommand{\SwiGLU}{\operatorname{SwiGLU}}
\newcommand{\SiLU}{\operatorname{SiLU}}
\newcommand{\Norm}{\operatorname{Norm}}
\newcommand{\Attn}{\operatorname{Attn}}
\newcommand{\tok}[1]{\textmd{\texttt{#1}}}
\newcommand{\pts}[1]{\textup{(#1~p)}}
\newcommand{\punto}{\textup{[1~p]}}
\newcommand{\apartado}[2]{\par\medskip\noindent\textbf{Apartado (#1)} \pts{#2}.\par\nopagebreak}
\newcommand{\repasar}[1]{\par\smallskip\noindent\textbf{Repasar:} #1.}

% La parte A, la parte B y el proyecto se numeran con letras.
\renewcommand{\thechapter}{\Alph{chapter}}
\setlocalecaption{spanish}{chapter}{Parte}
\crefname{chapter}{parte}{partes}

\hypersetup{pdftitle={Hito 3. Examen de hito: el Transformer},pdfauthor={Anatomía de un LLM}}

\begin{document}
\frontmatter
% ------------------------------------------------------------------ Portada
\begin{titlepage}
\centering
\vspace*{2.5cm}
{\large\scshape Anatomía de un LLM}\par
\vspace{0.4cm}
{\normalsize Fundamentos matemáticos de los modelos de lenguaje}\par
\vspace{2.2cm}
\rule{\textwidth}{0.8pt}\par\vspace{0.5cm}
{\small\scshape Fase III \textperiodcentered\ Hito 3 \textperiodcentered\ Examen de hito}\par\vspace{0.3cm}
{\Huge\bfseries Hito 3\par}\vspace{0.3cm}
{\LARGE El Transformer\par}
\vspace{0.3cm}
{\large Tokenización, atención, Transformer y mini-GPT\par}
\vspace{0.5cm}\rule{\textwidth}{0.4pt}\par
\vspace{1.6cm}
{\large Examen, proyecto de fase y soluciones}\par
\vfill
{\small Edición de octubre de 2026}\par
\end{titlepage}

% =====================================================================
\chapter*{Instrucciones}
\addcontentsline{toc}{chapter}{Instrucciones}
\markboth{Instrucciones}{}
% =====================================================================

Este examen cierra la fase~III del curso, dedicada al Transformer. Cubre los módulos~12 (tokenización y modelos secuenciales), 13 (atención), 14 (el Transformer decoder-only) y 15 (entrenamiento de un mini-GPT). Debe aprobarse antes de comenzar la fase~IV.

\section*{Condiciones}
\begin{center}
\begin{tabular}{@{}ll@{}}
\toprule
Duración & 2~h 30~min: unos 20~min para la parte~A y 2~h para la parte~B \\
Material permitido & calculadora científica y una hoja de fórmulas manuscrita (una cara) \\
Puntuación total & 100 puntos: parte~A, 20; parte~B, 80 \\
Aprobado & 70 puntos (70\,\%) \\
Módulos & M12, M13, M14 y M15 \\
\bottomrule
\end{tabular}
\end{center}

\section*{Estructura}
\begin{itemize}[leftmargin=1.4em,label=\textendash]
  \item \textbf{Parte A. Cuestiones} (20 puntos). Veinte preguntas de opción múltiple, de 1~punto cada una, con una sola respuesta correcta entre las cuatro propuestas. Las respuestas en blanco o incorrectas valen~0; no se penalizan los errores.
  \item \textbf{Parte B. Problemas} (80 puntos). Seis ejercicios de desarrollo, con la puntuación de cada apartado indicada entre paréntesis. Debe justificarse cada paso: un resultado sin razonamiento vale como máximo la mitad de su puntuación.
  \item \textbf{Proyecto de fase}. Se entrega aparte y se evalúa con la rúbrica del capítulo \hyperref[chap:proyecto]{«Proyecto de fase»}; no suma puntos al examen.
\end{itemize}

\section*{Convenios}
\begin{itemize}[leftmargin=1.4em,label=\textendash]
  \item Se redondea a tres decimales salvo que se indique otra cosa. Pueden usarse $e\approx 2{,}718$, $e^{2}\approx 7{,}389$ y $\ln 2\approx 0{,}693$.
  \item Las matrices de pesos multiplican por la derecha ($Y=XW$) y no hay sesgos salvo que se diga lo contrario.
  \item $1~\mathrm{GB}=10^{9}$~bytes. Una multiplicación seguida de una suma cuenta como 2~FLOPs.
  \item Los modelos con nombre inventado, como «Lince-1B», son ficticios. Los datos de GPT-2 y de LLaMA-2 son los publicados por sus autores.
  \item Las referencias del tipo «M13, §2.3» remiten al módulo y a la sección de los apuntes del curso.
\end{itemize}

\section*{Sobre las fuentes}
Los enunciados son de redacción propia. Los ejercicios marcados con \emph{Fuente} están adaptados de tareas y ejemplos de cursos universitarios publicados: la tarea~1 de Stanford CS336~\cite{cs336}, la tarea~3 de Stanford CS224N~\cite{cs224n} y los ejercicios del libro de Jurafsky y Martin~\cite{jm3}, texto de los cursos de procesamiento del lenguaje de Stanford. En todos los casos se han cambiado los datos para que el ejercicio pueda resolverse a mano.


\tableofcontents

\mainmatter
% =====================================================================
\chapter{Cuestiones}
\label{chap:parteA}
% =====================================================================

Cada cuestión vale 1~punto y tiene una sola respuesta correcta. Las cuestiones A.1 a A.5 corresponden al módulo~12; A.6 a A.10, al módulo~13; A.11 a A.18, al módulo~14, y A.19 y A.20, al módulo~15.

\begin{enumerate}[label=\textbf{A.\arabic*.},ref=A.\arabic*,leftmargin=2.6em,itemsep=8pt]

% ---------------------------------------------------------------- M12
\item\label{q:bpe-paso} ¿Qué hace el algoritmo BPE en cada paso de su entrenamiento? \punto
\begin{enumerate}[label=(\alph*)]
  \item Elimina del vocabulario el token menos probable.
  \item Divide la palabra más larga del corpus en dos mitades.
  \item Fusiona el par de símbolos adyacentes más frecuente en el corpus.
  \item Asigna a cada palabra un entero según su frecuencia.
\end{enumerate}

\item\label{q:bytes} Sin contar los tokens especiales, un tokenizador BPE a nivel de byte parte de un vocabulario de: \punto
\begin{enumerate}[label=(\alph*)]
  \item 128 símbolos.
  \item 256 símbolos.
  \item 65\,536 símbolos.
  \item tantos símbolos como caracteres distintos aparezcan en el corpus.
\end{enumerate}

\item\label{q:pinon} ¿Cuántos bytes ocupa la palabra «piñón» codificada en UTF-8? \punto
\begin{enumerate}[label=(\alph*)]
  \item 5.
  \item 6.
  \item 7.
  \item 10.
\end{enumerate}

\item\label{q:orden} Al codificar una palabra nueva con un tokenizador BPE ya entrenado, las reglas de fusión se aplican: \punto
\begin{enumerate}[label=(\alph*)]
  \item en el orden en que se aprendieron.
  \item de la que produce el token más largo a la que produce el más corto.
  \item en orden aleatorio, para dar variedad a la segmentación.
  \item solo la regla más frecuente.
\end{enumerate}

\item\label{q:rnn} ¿Por qué las redes recurrentes escalan peor que el Transformer? \punto
\begin{enumerate}[label=(\alph*)]
  \item Porque tienen más parámetros por capa.
  \item Porque no pueden usar embeddings.
  \item Porque solo admiten vocabularios de caracteres.
  \item Porque procesan los tokens uno tras otro, lo que impide paralelizar el entrenamiento a lo largo de la secuencia, y sus gradientes se desvanecen con la distancia.
\end{enumerate}

% ---------------------------------------------------------------- M13
\item\label{q:forma} En la atención de una cabeza sobre $n$ tokens, la matriz $A=\softmax\bigl(QK\T/\sqrt{d_k}\bigr)$ tiene forma: \punto
\begin{enumerate}[label=(\alph*)]
  \item $n\times d_k$.
  \item $d_k\times d_k$.
  \item $n\times n$.
  \item $n\times d_v$.
\end{enumerate}

\item\label{q:varianza} Si las componentes de $q,k\in\R^{d_k}$ son variables independientes de media~0 y varianza~1, la varianza de $q\cdot k$ es: \punto
\begin{enumerate}[label=(\alph*)]
  \item $1$.
  \item $d_k$.
  \item $\sqrt{d_k}$.
  \item $d_k^{2}$.
\end{enumerate}

\item\label{q:mascara} La máscara causal se implementa: \punto
\begin{enumerate}[label=(\alph*)]
  \item poniendo a~0 las puntuaciones de los tokens futuros antes de la softmax.
  \item poniendo a~0 los pesos de los tokens futuros después de la softmax, sin renormalizar.
  \item eliminando de $V$ las filas de los tokens futuros.
  \item sumando $-\infty$ a las puntuaciones de los tokens futuros antes de la softmax.
\end{enumerate}

\item\label{q:mha} Con $d=768$ y 12 cabezas, ¿cuántos parámetros tiene la atención multi-cabeza, sin sesgos? \punto
\begin{enumerate}[label=(\alph*)]
  \item 196\,608.
  \item 589\,824.
  \item 1\,179\,648.
  \item 2\,359\,296.
\end{enumerate}

\item\label{q:contexto} Pasar de un contexto de 1024 a uno de 8192 tokens multiplica la memoria de la matriz de atención, almacenada completa, por: \punto
\begin{enumerate}[label=(\alph*)]
  \item 8.
  \item 64.
  \item 512.
  \item 4096.
\end{enumerate}

% ---------------------------------------------------------------- M14
\item\label{q:prenorm} En un bloque pre-norm, la subcapa de atención se escribe: \punto
\begin{enumerate}[label=(\alph*)]
  \item $h=x+\Attn(\Norm(x))$.
  \item $h=\Norm(x+\Attn(x))$.
  \item $h=\Norm(\Attn(x))$.
  \item $h=\Attn(x)+\Norm(x)$.
\end{enumerate}

\item\label{q:normfinal} ¿Por qué un Transformer pre-norm necesita una normalización final antes del unembedding? \punto
\begin{enumerate}[label=(\alph*)]
  \item Porque la norma del flujo residual crece con la profundidad.
  \item Porque la softmax exige entradas positivas.
  \item Para poder atar los pesos de entrada y de salida.
  \item No la necesita; solo la llevan los modelos post-norm.
\end{enumerate}

\item\label{q:swiglu} Un MLP SwiGLU tiene el mismo número de parámetros que un MLP clásico de dimensión oculta $4d$ cuando su dimensión oculta vale: \punto
\begin{enumerate}[label=(\alph*)]
  \item $2d$.
  \item $\tfrac{8}{3}d$.
  \item $3d$.
  \item $4d$.
\end{enumerate}

\item\label{q:rope-rel} RoPE consigue que la puntuación $\tilde q_m\cdot\tilde k_n$ entre el token~$m$ y el token~$n$: \punto
\begin{enumerate}[label=(\alph*)]
  \item dependa de las posiciones solo a través de $n-m$.
  \item dependa de las posiciones solo a través de $m+n$.
  \item sea independiente de las posiciones.
  \item crezca linealmente con la posición.
\end{enumerate}

\item\label{q:rope-v} RoPE gira las consultas y las claves, pero no los valores, porque: \punto
\begin{enumerate}[label=(\alph*)]
  \item $v$ tiene dimensión impar.
  \item girar $v$ cambiaría su norma.
  \item la posición debe influir en a quién se atiende, no en qué contenido se copia.
  \item $v$ no participa en la atención.
\end{enumerate}

\item\label{q:llama-emb} LLaMA-2~7B no ata sus embeddings ($V=32\,000$, $d=4096$). ¿Cuántos parámetros suman la matriz de entrada y la de salida? \punto
\begin{enumerate}[label=(\alph*)]
  \item 65\,536\,000.
  \item 131\,072\,000.
  \item 262\,144\,000.
  \item 524\,288\,000.
\end{enumerate}

\item\label{q:16bytes} Entrenar con Adam en precisión mixta (pesos y gradientes en bf16, copia maestra en fp32 y momentos $m$ y $v$ en fp32) ocupa, sin contar las activaciones: \punto
\begin{enumerate}[label=(\alph*)]
  \item 4 bytes por parámetro.
  \item 8 bytes por parámetro.
  \item 12 bytes por parámetro.
  \item 16 bytes por parámetro.
\end{enumerate}

\item\label{q:6nd} Según la regla $6ND$, ¿cuántos FLOPs cuesta, aproximadamente, entrenar un modelo de $10^{9}$ parámetros con $2\cdot 10^{10}$ tokens? \punto
\begin{enumerate}[label=(\alph*)]
  \item $6\cdot 10^{18}$.
  \item $2\cdot 10^{19}$.
  \item $4\cdot 10^{19}$.
  \item $1{,}2\cdot 10^{20}$.
\end{enumerate}

% ---------------------------------------------------------------- M15
\item\label{q:lnV} Un GPT con $V=50\,257$ tokens, recién inicializado con logits casi nulos, debe mostrar una pérdida inicial (en nats por token) cercana a: \punto
\begin{enumerate}[label=(\alph*)]
  \item $10{,}82$.
  \item $1$.
  \item $0$.
  \item $50\,257$.
\end{enumerate}

\item\label{q:checkpoint} Durante el entrenamiento, la pérdida de entrenamiento sigue bajando y la de validación empieza a subir. ¿Qué punto de control debe conservarse como modelo final? \punto
\begin{enumerate}[label=(\alph*)]
  \item El último.
  \item El de menor pérdida de validación.
  \item El de menor pérdida de entrenamiento.
  \item El primero, porque aún no ha sobreajustado.
\end{enumerate}

\end{enumerate}

% =====================================================================
\chapter{Problemas}
\label{chap:parteB}
% =====================================================================

La parte~B consta de seis ejercicios con un total de 80~puntos. La puntuación de cada apartado figura entre paréntesis. Debe justificarse cada paso y expresarse cada resultado con sus unidades.

% ---------------------------------------------------------------------
\section*{Módulo 14. El bloque decoder}

\begin{ejercicio}{2}{eje:B-bloque}
\textbf{(15 puntos.)} Un Transformer decoder-only del estilo de LLaMA recibe $n$ tokens y tiene dimensión de modelo $d$, $L$ bloques, $h$ cabezas de atención, MLP SwiGLU de dimensión oculta $\dff$, normalización RMSNorm, codificación de posición RoPE y un vocabulario de $V$ tokens con embeddings sin atar.
\begin{enumerate}[label=(\alph*)]
  \item \pts{5} Dibuje el diagrama completo, desde los tokens enteros hasta la distribución del token siguiente, con un bloque detallado. Señale las dos conexiones residuales y las dos normalizaciones del bloque, el lugar en que entra la información de posición y la normalización final.
  \item \pts{3} Escriba las ecuaciones del bloque $\ell$, con entrada $h^{\ell-1}$ y salida $h^{\ell}$, y la ecuación de los logits.
  \item \pts{3} Indique la forma de: la entrada al primer bloque; $Q$ tras separar las cabezas; la matriz de atención de una cabeza; la salida de la proyección $W_{\mathrm{gate}}$; la salida de un bloque, y los logits.
  \item \pts{2} ¿Qué subcapa mezcla información \emph{entre} tokens y cuál procesa cada token \emph{por separado}? ¿Por qué cada subcapa debe devolver un vector de $\R^{d}$?
  \item \pts{2} Dé dos razones por las que los modelos de lenguaje actuales prefieren la variante pre-norm a la post-norm.
\end{enumerate}
\fuente{adaptado de Stanford CS336~\cite{cs336}, tarea~1, §3.4 y §3.5, ecuación~(15) y problemas \emph{transformer\_block} y \emph{transformer\_lm}, y §7.3, problema \emph{pre\_norm\_ablation} (la implementación y el experimento se sustituyen por el diagrama, las ecuaciones y el análisis).}
\end{ejercicio}

% ---------------------------------------------------------------------
\section*{Módulo 13. Atención con máscara causal}

\begin{ejercicio}{2}{eje:B-atencion}
\textbf{(15 puntos.)} Tres tokens tienen embeddings de dimensión $d=2$: $x_1=(1,0)$, $x_2=(0,1)$ y $x_3=(1,1)$, apilados como filas de una matriz $X\in\R^{3\times 2}$. Una cabeza de atención usa $d_k=4$, $d_v=2$ y
\[
W_Q=\begin{pmatrix}1&1&0&0\\0&0&1&1\end{pmatrix},\qquad
W_K=\begin{pmatrix}2&0&0&0\\0&0&2&0\end{pmatrix},\qquad
W_V=\begin{pmatrix}1&1\\0&1\end{pmatrix}.
\]
\begin{enumerate}[label=(\alph*)]
  \item \pts{3} Calcule $Q=XW_Q$, $K=XW_K$ y $V=XW_V$.
  \item \pts{3} Calcule $S=QK\T$ y $S/\sqrt{d_k}$.
  \item \pts{2} Escriba la matriz de máscara causal $M$ y la matriz $Z=S/\sqrt{d_k}+M$. ¿Por qué la máscara usa $-\infty$ y no~$0$?
  \item \pts{3} Calcule $A=\softmax(Z)$, aplicada por filas, y la salida $O=AV$.
  \item \pts{2} Se sustituye $x_3$ por $(5,-3)$. Sin hacer cuentas largas, indique qué filas de $O$ cambian, justificándolo con la estructura de $A$, y calcule la nueva fila $o_3$. ¿Qué hace la cabeza en este caso?
  \item \pts{2} Calcule $\partial a_{22}/\partial z_{22}$ con los valores del apartado~(d). ¿Qué ocurriría con esta derivada si no se dividiera por $\sqrt{d_k}$ y las puntuaciones fueran muy grandes?
\end{enumerate}
\fuente{adaptado de Stanford CS224N~\cite{cs224n}, tarea~3, pregunta~1(a), «Copying in attention» (apartado~e), y de Jurafsky y Martin~\cite{jm3}, cap.~7, ejercicio~7.2 (apartado~c); datos numéricos propios.}
\end{ejercicio}

% ---------------------------------------------------------------------
\section*{Módulo 12. BPE a mano}

\begin{ejercicio}{2}{eje:B-bpe}
\textbf{(12 puntos.)} Un corpus contiene estas palabras con sus frecuencias: \tok{pino}~$\times 4$, \tok{pinos}~$\times 2$, \tok{vino}~$\times 3$, \tok{vinos}~$\times 1$ y \tok{pinta}~$\times 2$. Los símbolos iniciales son los caracteres y las fusiones no cruzan los límites entre palabras.
\begin{enumerate}[label=(\alph*)]
  \item \pts{1} ¿Cuántos símbolos tiene el corpus al empezar?
  \item \pts{6} Aplique cuatro fusiones BPE. En cada paso indique el par elegido y su frecuencia, el par que queda en segundo lugar y el número total de símbolos del corpus tras la fusión.
  \item \pts{3} Codifique con las cuatro reglas obtenidas las palabras \tok{vinos}, \tok{pinto}, \tok{sino} y \tok{ovino}, e indique el número de tokens de cada una.
  \item \pts{2} Un tokenizador BPE a nivel de byte procesa la palabra «piñón». ¿Cuántos símbolos iniciales tiene? ¿Qué pares de bytes conviene fusionar pronto en un corpus en español y por qué?
\end{enumerate}
\fuente{adaptado de Stanford CS336~\cite{cs336}, tarea~1, §2.4, ejemplo \emph{bpe\_example}, y §2.6.1, ejemplo \emph{bpe\_encoding}, y de Jurafsky y Martin~\cite{jm3}, §2.4.1 y ejercicio~2.3 (corpus y palabras cambiados).}
\end{ejercicio}

% ---------------------------------------------------------------------
\section*{Módulo 14. Parámetros, FLOPs y memoria}

\begin{ejercicio}{3}{eje:B-lince}
\textbf{(15 puntos.)} El modelo ficticio \textbf{Lince-1B} tiene $d=2048$, $L=24$, $h=16$, MLP SwiGLU, $V=48\,000$, embeddings \emph{atados}, RMSNorm con solo vector de pesos, RoPE y ningún sesgo.
\begin{enumerate}[label=(\alph*)]
  \item \pts{2} Calcule $\tfrac{8}{3}d$ y redondéelo al múltiplo de 256 inmediatamente superior para obtener $\dff$.
  \item \pts{4} Calcule el número de parámetros de un bloque y el total~$N$. Compare $N$ con la aproximación $12Ld^{2}$ e indique qué fracción de $N$ corresponde a los MLP.
  \item \pts{3} Estime los FLOPs por token en la generación y los FLOPs del entrenamiento con $D=2\cdot 10^{11}$ tokens. ¿Cuántas horas de GPU H100 ($989\cdot 10^{12}$~FLOP/s en bf16) se necesitan con un aprovechamiento del 40\,\%? ¿Cuántas horas de reloj con 64~GPU?
  \item \pts{3} Calcule la memoria que ocupan los pesos en bf16 y en int4, y la del estado del entrenamiento con Adam en precisión mixta. ¿Cabe el entrenamiento en una GPU de 24~GB?
  \item \pts{3} La atención añade unos $4Lnd$ FLOPs por token con un contexto de $n$ tokens. ¿Qué fracción de $2N$ supone con $n=4096$? ¿Para qué $n$ iguala a $2N$?
\end{enumerate}
\fuente{adaptado de Stanford CS336~\cite{cs336}, tarea~1, §3.5, problema \emph{transformer\_accounting}, apartados (a) a~(e), y §4.3, problema \emph{adamw\_accounting}, apartado~(d), y de Jurafsky y Martin~\cite{jm3}, cap.~7, ejercicio~7.1 (se sustituye GPT-2~XL por un modelo ficticio y la precisión fp32 por la mixta).}
\end{ejercicio}

% ---------------------------------------------------------------------
\section*{Módulos 13 y 14. Posiciones y coste de la atención}

\begin{ejercicio}{2}{eje:B-rope}
\textbf{(11 puntos.)}
\begin{enumerate}[label=(\alph*)]
  \item \pts{5} Se aplica RoPE en dimensión $d_k=2$ con $\theta=\pi/4$, $q=(1,0)$ y $k=(1,1)$. Use la identidad
  \[
  q\T R(\varphi)\,k=(q\cdot k)\cos\varphi-(q_1k_2-q_2k_1)\sin\varphi,\qquad \varphi=(n-m)\theta,
  \]
  para calcular la puntuación entre la consulta en la posición~$m$ y la clave en la posición~$n$ en los pares $(m,n)=(2,1)$, $(5,4)$, $(3,3)$, $(1,2)$ y $(0,4)$. ¿Qué se observa en los dos primeros?
  \item \pts{2} En Lince-1B (\cref{eje:B-lince}), ¿cuánto vale $d_k$? Si se duplica $h$ manteniendo $d$, ¿cuántos parámetros de atención se ganan?
  \item \pts{4} Con un contexto de $n=32\,768$ tokens, calcule la memoria en fp16 de la matriz de atención de una cabeza y la de todas las cabezas y capas de Lince-1B. ¿A partir de qué $n$ dominan los FLOPs cuadráticos de la atención, $4n^{2}d$ por capa, sobre los de las proyecciones, $8nd^{2}$ por capa? ¿Qué técnica evita almacenar $A$ completa?
\end{enumerate}
\fuente{adaptado de Stanford CS336~\cite{cs336}, tarea~1, §3.4.3, ecuaciones~(8) y~(9) (RoPE), y §3.5, problema \emph{transformer\_accounting}, apartado~(e) (se pide el cálculo a mano en dimensión~2 y la memoria de $A$).}
\end{ejercicio}

% ---------------------------------------------------------------------
\section*{Módulo 15. Entrenar y leer un mini-GPT}

\begin{ejercicio}{2}{eje:B-minigpt}
\textbf{(12 puntos.)} Un mini-GPT de caracteres tiene $V=80$, contexto $T=128$, $d=128$, $L=4$, $h=4$, MLP clásico de dimensión oculta $4d$ con GELU, LayerNorm con peso y sesgo, posiciones aprendidas, embeddings atados y ninguna capa lineal con sesgo. Se entrena con lotes de $B=32$ secuencias durante 5000~pasos sobre un corpus de $10^{6}$ caracteres. El \cref{tab:curvas} recoge sus pérdidas medias.

\begin{table}[H]
\centering\small
\setlength{\tabcolsep}{3.6pt}
\caption{Pérdidas del mini-GPT, en nats por carácter.}\label{tab:curvas}
\begin{tabular}{@{}lccccccccccc@{}}
\toprule
Paso & 0 & 500 & 1000 & 1500 & 2000 & 2500 & 3000 & 3500 & 4000 & 4500 & 5000 \\
\midrule
Entrenamiento & 4{,}39 & 2{,}12 & 1{,}66 & 1{,}47 & 1{,}38 & 1{,}30 & 1{,}22 & 1{,}15 & 1{,}08 & 1{,}03 & 0{,}98 \\
Validación & 4{,}39 & 2{,}14 & 1{,}70 & 1{,}52 & 1{,}46 & 1{,}44 & 1{,}45 & 1{,}48 & 1{,}52 & 1{,}57 & 1{,}62 \\
\bottomrule
\end{tabular}
\end{table}

\begin{enumerate}[label=(\alph*)]
  \item \pts{3} Cuente los parámetros del modelo.
  \item \pts{2} Explique por qué la pérdida inicial vale unos $4{,}39$ nats.
  \item \pts{3} ¿Qué punto de control debe conservarse? Calcule su perplejidad de validación y sus bits por carácter. Un modelo de bigramas alcanza una pérdida de validación de $2{,}45$ nats: ¿cuál es su perplejidad?
  \item \pts{2} ¿Cuántos tokens procesa el entrenamiento? ¿A cuántas épocas equivale? Relacione el resultado con el \cref{tab:curvas} y proponga dos remedios.
  \item \pts{2} Al generar, el modelo produce los logits $(2,1,0)$ para tres caracteres. Calcule las probabilidades con temperatura $\tau=1$ y $\tau=0{,}5$. ¿Qué efecto tiene bajar la temperatura sobre las muestras?
\end{enumerate}
\fuente{adaptado de Stanford CS336~\cite{cs336}, tarea~1, §4.1, ecuación~(18) (perplejidad), y §6, ecuación~(23) y problema \emph{decoding} (temperatura), y de Jurafsky y Martin~\cite{jm3}, cap.~7, ejercicios~7.4 y~7.5 (curvas y logits propios).}
\end{ejercicio}

% =====================================================================
\chapter{Proyecto de fase}
\label{chap:proyecto}
% =====================================================================

\section*{Mini-GPT en español con informe técnico}

\paragraph{Objetivo.} Entrenar desde cero, en PyTorch y en CPU, un Transformer decoder-only de caracteres, y redactar un informe técnico con las curvas de pérdida, muestras generadas e interpretación de mapas de atención. El proyecto integra los cuatro módulos de la fase: el vocabulario de caracteres (M12), la atención causal multi-cabeza (M13), el bloque pre-norm (M14) y el bucle de entrenamiento y generación (M15).

\paragraph{Datos.} Se usará un corpus en español escrito por el estudiante o generado con plantillas (frases que combinan al azar sujetos, verbos, objetos y lugares). No deben usarse textos literarios existentes, ni siquiera de dominio público. Bastan entre $300\,000$ y $10^{6}$ caracteres.

\section*{Entregables}
\begin{enumerate}[label=\arabic*.]
  \item \texttt{H3-minigpt.ipynb}: código completo, que se ejecute en CPU en menos de 10~minutos.
  \item \texttt{H3-informe.md}, de 2 a 4 páginas, con:
  \begin{itemize}[label=\textendash]
    \item la configuración ($d$, $L$, $h$, $T$, $V$, $B$, número de pasos y tasa de aprendizaje) y el número de parámetros, contado a mano y con código;
    \item las curvas de pérdida de entrenamiento y de validación en la misma gráfica, con la pérdida de un modelo de bigramas como línea de referencia y el valor $\ln V$ marcado en el paso~0;
    \item tres muestras de 300~caracteres con temperaturas $0{,}5$, $1{,}0$ y $1{,}5$, comentadas;
    \item mapas de atención de al menos dos cabezas en dos capas sobre una frase del corpus, con su interpretación (qué cabeza atiende al carácter anterior, al inicio de palabra o a sí misma);
    \item un experimento de ablación: eliminar las posiciones, usar una cabeza en lugar de cuatro o pasar de pre-norm a post-norm, comparando las pérdidas.
  \end{itemize}
\end{enumerate}

\section*{Pasos sugeridos}
\begin{enumerate}[label=\arabic*.]
  \item Generar el corpus y dividirlo en un 90\,\% para entrenamiento y un 10\,\% para validación, con un corte contiguo, no aleatorio.
  \item Calcular la pérdida de validación de un modelo de bigramas con suavizado de Laplace (módulo~1): es el listón que hay que superar.
  \item Implementar el bloque pre-norm con atención causal multi-cabeza y MLP. Inicializar los embeddings con desviación típica $0{,}02$: con $\mathcal N(0,1)$ y pesos atados, la pérdida inicial queda muy por encima de $\ln V$ (M15, §1.4).
  \item Comprobar la causalidad: al cambiar el último carácter, los logits de las posiciones anteriores no deben cambiar (M13, §3.2).
  \item Entrenar con AdamW, calentamiento lineal y decaimiento coseno, y recorte del gradiente a norma~$1{,}0$ (M15, §3.2 a §3.4). Evaluar cada 100~pasos y conservar el punto de control con menor pérdida de validación (M15, §4.2).
  \item Generar muestras y extraer las matrices $A$ de cada capa, devueltas por \texttt{forward}. Dibujarlas con \texttt{imshow}, con los caracteres en los ejes.
\end{enumerate}

\paragraph{Orientación.} En una prueba con un corpus de plantillas de $717\,000$ caracteres ($V=33$), un modelo con $d=64$, $L=2$, $h=4$ y $T=32$ ($103\,104$ parámetros) redujo la pérdida de validación de $3{,}50$ a $0{,}36$ nats en 800~pasos, frente a $1{,}75$ del modelo de bigramas, y tardó menos de 30~s en CPU. Las muestras eran ya frases gramaticales del estilo de las plantillas.

\section*{Criterios de evaluación}
El proyecto se califica sobre 100 con los pesos del \cref{tab:rubrica}. Se considera superado con 70 o más.

\begin{table}[ht]
\centering\footnotesize
\renewcommand{\arraystretch}{1.15}
\caption{Rúbrica del proyecto de fase.}\label{tab:rubrica}
\begin{tabularx}{\textwidth}{@{}>{\raggedright\arraybackslash}p{2.45cm}c>{\raggedright\arraybackslash}X>{\raggedright\arraybackslash}X>{\raggedright\arraybackslash}X@{}}
\toprule
Criterio & Peso & Excelente & Suficiente & Insuficiente \\
\midrule
Implementación del modelo & 25\,\% & Bloque pre-norm correcto, causalidad comprobada con una prueba, recuento a mano igual al del código & Funciona y entrena; recuento con pequeñas discrepancias explicadas & Errores en la máscara o en las formas; no entrena \\
\addlinespace
Entrenamiento y curvas & 20\,\% & Curvas de entrenamiento y validación con bigramas y $\ln V$ como referencias; validación $\le 75\,\%$ de la de bigramas; mejor punto de control justificado & Curvas correctas y validación por debajo de bigramas & Solo pérdida de entrenamiento, o no supera a bigramas \\
\addlinespace
Muestras generadas & 15\,\% & Tres temperaturas comentadas con criterio (repetición frente a incoherencia) & Muestras presentadas con un comentario breve & Sin muestras o sin comentar \\
\addlinespace
Mapas de atención & 20\,\% & Mapas legibles con ejes etiquetados, triángulo inferior visible, interpretación de cada cabeza con ejemplos & Mapas correctos con interpretación genérica & Mapas sin máscara visible o sin interpretar \\
\addlinespace
Ablación & 10\,\% & Experimento controlado (misma semilla y número de pasos) con conclusión cuantitativa & Experimento realizado con comparación cualitativa & No hay ablación \\
\addlinespace
Informe & 10\,\% & Claro y reproducible, con cifras en formato español y relación con los LLM reales & Completo pero desordenado & Faltan secciones \\
\bottomrule
\end{tabularx}
\end{table}


\appendix
% =====================================================================
\chapter{Soluciones}
\label{ap:soluciones}
% =====================================================================

Las soluciones de la parte~B se presentan apartado por apartado, en cinco partes: datos, planteamiento, desarrollo, resultado y comprobación. Se indican también los criterios de corrección cuando el reparto de puntos de un apartado no es evidente. Todos los valores numéricos se han comprobado con Python; en los pasos intermedios se conservan cuatro cifras significativas.

% ---------------------------------------------------------------------
\section{Clave de respuestas de la parte A}

\begin{table}[H]
\centering\footnotesize
\renewcommand{\arraystretch}{1.05}
\caption{Clave de respuestas de la parte~A.}\label{tab:clave}
\begin{tabularx}{\textwidth}{@{}clX>{\raggedright\arraybackslash}p{1.95cm}@{}}
\toprule
N.º & Resp. & Justificación & Apuntes \\
\midrule
A.1 & (c) & BPE cuenta los pares adyacentes, ponderados por frecuencia, y fusiona el más frecuente. & M12, §2.1 \\
A.2 & (b) & Parte de los 256 valores posibles de un byte; por eso nunca aparece un token desconocido. & M12, §2.4 \\
A.3 & (c) & «p», «i» y «n» (ASCII) ocupan 1~byte cada una; «ñ» y «ó», 2 cada una: $3\cdot 1+2\cdot 2=7$. & M12, §1.3 \\
A.4 & (a) & Las reglas se aplican en el orden en que se aprendieron; así la codificación es determinista. & M12, §2.2 \\
A.5 & (d) & La recurrencia $h_t=f(h_{t-1},x_t)$ exige $n$ pasos secuenciales y el gradiente se multiplica $n$ veces. & M12, §5.5 \\
A.6 & (c) & Cada token (fila) tiene un peso para cada token (columna). & M13, §2.1 \\
A.7 & (b) & $q\cdot k$ es la suma de $d_k$ productos independientes de media~0 y varianza~1. & M13, §2.3 \\
A.8 & (d) & $e^{-\infty}=0$ y los pesos permitidos siguen sumando~1; un~0 antes de la softmax no anula nada. & M13, §3.1 \\
A.9 & (d) & $4d^{2}=4\cdot 768^{2}=2\,359\,296$, sin dependencia de $h$; (c) es $2d^{2}$. & M13, §4.2 \\
A.10 & (b) & La memoria de $A$ crece como $n^{2}$: $(8192/1024)^{2}=8^{2}=64$. & M13, §4.4 \\
A.11 & (a) & Pre-norm normaliza la entrada de la subcapa y deja limpio el camino residual; (b) es post-norm. & M14, §2.2 \\
A.12 & (a) & En pre-norm, $h^{L}=h^{0}+\sum$ aportaciones y su norma crece con $L$. & M14, §2.3 \\
A.13 & (b) & $3d\,\dff=2d\cdot 4d=8d^{2}$, luego $\dff=\tfrac{8}{3}d$. & M14, §3.3 \\
A.14 & (a) & $R_m\T R_n=R_{n-m}$, luego $\tilde q_m\cdot\tilde k_n=q\T R_{n-m}k$. & M14, §4.3 \\
A.15 & (c) & $v$ es el contenido que se copia; la posición solo debe decidir los pesos. & M14, §4.3 \\
A.16 & (c) & $2Vd=2\cdot 32\,000\cdot 4096=262\,144\,000$. & M14, §1.4 y §5.2 \\
A.17 & (d) & $2+2+4+4+4=16$~bytes por parámetro. & M14, §5.3 \\
A.18 & (d) & $6ND=6\cdot 10^{9}\cdot 2\cdot 10^{10}=1{,}2\cdot 10^{20}$. & M14, §6.3 \\
A.19 & (a) & Logits iguales dan la distribución uniforme y la pérdida $\ln V=\ln 50\,257=10{,}82$. & M15, §1.4 \\
A.20 & (b) & La pérdida de validación estima el rendimiento sobre datos nuevos. & M15, §4.2 y §4.4 \\
\bottomrule
\end{tabularx}
\end{table}

% ---------------------------------------------------------------------
\section{Soluciones de la parte B}

% ===================================================================== B1
\begin{solucion}{eje:B-bloque}

\apartado{a}{5}
\datos Transformer decoder-only del estilo de LLaMA: $n$ tokens, dimensión $d$, $L$ bloques, $h$ cabezas, SwiGLU con $\dff$, RMSNorm, RoPE y $V$ tokens sin atar.
\planteamiento recorrer el camino de los datos: enteros, embedding, $L$ bloques pre-norm, normalización final, unembedding y softmax (M14, §1.2).
\pasos
\begin{enumerate}[label=\arabic*.]
  \item Entrada: los $n$ enteros seleccionan filas de $W_E\in\R^{V\times d}$ y forman $h^{0}\in\R^{n\times d}$.
  \item Primera subcapa del bloque: RMSNorm~1, proyecciones $Q$, $K$, $V$, rotación RoPE de $q$ y $k$ (aquí entra la posición), atención causal en cada cabeza, concatenación y $W_O$, y suma residual~1.
  \item Segunda subcapa: RMSNorm~2, SwiGLU con $W_{\mathrm{gate}}$, $W_{\mathrm{up}}$ y $W_{\mathrm{down}}$, y suma residual~2.
  \item Tras los $L$ bloques: RMSNorm final, $W_U\in\R^{d\times V}$ y softmax.
\end{enumerate}
\resultado \textbf{el diagrama de la \cref{fig:bloque}}, con dos residuales por bloque, una normalización antes de cada subcapa, RoPE aplicada a $q$ y $k$ y una RMSNorm final.
\comprobacion todas las flechas entre bloques transportan matrices $n\times d$; solo cambian de forma la entrada (enteros a $\R^{d}$) y la salida ($\R^{d}$ a $\R^{V}$).

\begin{figure}[ht]
\centering
\begin{tikzpicture}[
  caja/.style={draw,rounded corners=2pt,minimum width=7.2cm,minimum height=0.62cm,font=\small,align=center,fill=white},
  suma/.style={draw,circle,inner sep=0.5pt,font=\small,fill=white},
  flecha/.style={-{Stealth[length=2.2mm]}},
  node distance=0.36cm]
\node[caja] (tok) {Tokens $x_1,\dots,x_n$ (enteros)};
\node[caja,below=of tok] (emb) {Embedding $W_E$ ($V\times d$) $\;\to\; h^{0}$ ($n\times d$)};
\node[caja,below=0.75cm of emb] (n1) {RMSNorm 1};
\node[caja,below=of n1] (qkv) {$Q=XW_Q$, $K=XW_K$, $V=XW_V$};
\node[caja,below=of qkv] (rope) {RoPE sobre $q$ y $k$ (entra la posición)};
\node[caja,below=of rope] (att) {$h$ cabezas: $\softmax\bigl(Q_iK_i\T/\sqrt{d_k}+M\bigr)V_i$};
\node[caja,below=of att] (wo) {Concatenar y $W_O$ ($d\times d$)};
\node[suma,below=of wo] (s1) {$+$};
\node[caja,below=of s1] (n2) {RMSNorm 2};
\node[caja,below=of n2] (mlp) {SwiGLU: $\bigl(\SiLU(XW_{\mathrm{gate}})\odot XW_{\mathrm{up}}\bigr)W_{\mathrm{down}}$};
\node[suma,below=of mlp] (s2) {$+$};
\node[caja,below=0.75cm of s2] (nf) {RMSNorm final};
\node[caja,below=of nf] (wu) {Unembedding $W_U$ ($d\times V$): logits $n\times V$};
\node[caja,below=of wu] (sm) {softmax $\to P(x_{n+1}\mid x_{1:n})$};
\foreach \a/\b in {tok/emb,emb/n1,n1/qkv,qkv/rope,rope/att,att/wo,wo/s1,s1/n2,n2/mlp,mlp/s2,s2/nf,nf/wu,wu/sm}
  \draw[flecha] (\a) -- (\b);
\coordinate (r0) at ($(emb.south)!0.5!(n1.north)$);
\draw[flecha] (r0) -- ++(-4.3,0) |- (s1.west) node[pos=0.25,left,font=\footnotesize,align=right] {residual\\1};
\coordinate (r1) at ($(s1.south)!0.5!(n2.north)$);
\draw[flecha] (r1) -- ++(-4.3,0) |- (s2.west) node[pos=0.25,left,font=\footnotesize,align=right] {residual\\2};
\node[draw,dashed,rounded corners=4pt,inner sep=5pt,fit=(n1)(qkv)(rope)(att)(wo)(s1)(n2)(mlp)(s2)] (bl) {};
\node[anchor=north west,font=\footnotesize\itshape] at (bl.north east) {$\times L$};
\end{tikzpicture}
\caption{Transformer decoder-only con bloques pre-norm del estilo de LLaMA. El recuadro discontinuo es un bloque y se repite $L$ veces.}\label{fig:bloque}
\end{figure}

Criterios: embedding y unembedding con sus formas (1~p); dos residuales que puentean cada subcapa (1~p); normalizaciones antes de cada subcapa y normalización final (1~p); atención causal con máscara y RoPE en $q$ y $k$ (1~p); MLP con sus tres matrices (1~p). Si se dibujan posiciones aprendidas sumadas al embedding, al estilo de GPT-2, se acepta siempre que se indique.

\apartado{b}{3}
\datos entrada del bloque $h^{\ell-1}$; arquitectura pre-norm.
\planteamiento en pre-norm, cada subcapa lee su entrada normalizada y suma su salida al flujo residual (M14, §2.2).
\pasos
\begin{enumerate}[label=\arabic*.]
  \item Atención: $a^{\ell}=h^{\ell-1}+\MHA\bigl(\RMSNorm_1(h^{\ell-1})\bigr)$.
  \item MLP: $h^{\ell}=a^{\ell}+\SwiGLU\bigl(\RMSNorm_2(a^{\ell})\bigr)$.
  \item Salida: $\mathrm{logits}=\RMSNorm_f(h^{L})\,W_U$.
\end{enumerate}
\resultado \textbf{las tres ecuaciones anteriores}, 1~p cada una.
\comprobacion desarrollando la recurrencia, $h^{L}=h^{0}+\sum_{\ell}\bigl(\text{aportación de la atención}+\text{aportación del MLP}\bigr)$: el camino identidad no atraviesa ninguna normalización (M14, §2.1).

\apartado{c}{3}
\datos $n$ tokens, $d$, $h$ cabezas, $\dff$, $V$.
\planteamiento seguir las formas en cada producto $Y=XW$; 0,5~p por forma.
\pasos
\begin{enumerate}[label=\arabic*.]
  \item Entrada al primer bloque: $n$ filas de $W_E$, es decir, $n\times d$.
  \item $Q=XW_Q$ es $n\times d$; separada en $h$ cabezas de dimensión $d/h$: $h\times n\times (d/h)$.
  \item Atención de una cabeza: $(n\times d/h)(d/h\times n)$, es decir, $n\times n$.
  \item $W_{\mathrm{gate}}\in\R^{d\times\dff}$: $(n\times d)(d\times\dff)$ da $n\times\dff$.
  \item Salida de un bloque: $n\times d$, la misma forma que la entrada.
  \item Logits: $(n\times d)(d\times V)$ da $n\times V$.
\end{enumerate}
\resultado \textbf{$n\times d$; $h\times n\times(d/h)$; $n\times n$; $n\times\dff$; $n\times d$; $n\times V$.}
\comprobacion como $h\cdot(d/h)=d$, la concatenación de las cabezas recupera la forma $n\times d$.

\apartado{d}{2}
\datos subcapas de atención y MLP.
\planteamiento comprobar si la operación combina filas distintas de la matriz $n\times d$ (M14, §3.1).
\pasos
\begin{enumerate}[label=\arabic*.]
  \item La atención mezcla tokens: la fila $i$ de $O=AV$ es $\sum_{j\le i}a_{ij}v_j$.
  \item El MLP aplica la misma función a cada fila por separado: la salida del token~$i$ solo depende de la fila~$i$.
  \item Cada subcapa \emph{suma} su salida al flujo residual, que vive en $\R^{d}$; la suma exige la misma dimensión, y por eso $W_O\in\R^{d\times d}$ y $W_{\mathrm{down}}\in\R^{\dff\times d}$ devuelven a $\R^{d}$.
\end{enumerate}
\resultado \textbf{la atención mezcla información entre tokens; el MLP actúa token a token; ambas devuelven $\R^{d}$ porque se suman al flujo residual.}
\comprobacion si $W_{\mathrm{down}}$ fuera $\dff\times\dff$, la suma $a^{\ell}+\SwiGLU(\cdot)$ no estaría definida.

\apartado{e}{2}
\datos pre-norm, $h=x+F(\Norm(x))$, frente a post-norm, $h=\Norm(x+F(x))$.
\planteamiento comparar el camino del gradiente y la estabilidad del entrenamiento (M14, §2.2).
\pasos
\begin{enumerate}[label=\arabic*.]
  \item En pre-norm existe un camino identidad de la entrada a la salida: el jacobiano $\partial h^{L}/\partial h^{0}$ contiene el término $I$ y el gradiente llega a las primeras capas sin atravesar normalizaciones.
  \item El entrenamiento es estable con muchas capas y con un calentamiento corto o nulo; en post-norm, el gradiente de las últimas capas es grande en la inicialización y obliga a un calentamiento largo y cuidadoso~\cite{xiong}.
\end{enumerate}
\resultado \textbf{gradiente limpio por el camino identidad y entrenamiento estable en profundidad, poco sensible al calentamiento}, 1~p por razón.
\comprobacion GPT-2~\cite{gpt2}, LLaMA-2~\cite{llama2} y la tarea~1 de CS336~\cite[§3.4]{cs336} usan pre-norm; el problema \emph{pre\_norm\_ablation} de esa tarea compara ambas variantes experimentalmente.

\repasar{M14, §1.2, §2.1, §2.2, §2.3 y §3.1}
\end{solucion}

% ===================================================================== B2
\begin{solucion}{eje:B-atencion}

\apartado{a}{3}
\datos $X$ con filas $x_1=(1,0)$, $x_2=(0,1)$, $x_3=(1,1)$; $W_Q,W_K\in\R^{2\times 4}$, $W_V\in\R^{2\times 2}$.
\planteamiento la fila $i$ de $XW$ es $x_iW$: $x_1$ selecciona la primera fila de $W$, $x_2$ la segunda y $x_3$ suma ambas (M13, §2.1).
\pasos
\begin{enumerate}[label=\arabic*.]
  \item $Q$: $(1,1,0,0)$, $(0,0,1,1)$ y su suma $(1,1,1,1)$.
  \item $K$: $(2,0,0,0)$, $(0,0,2,0)$ y su suma $(2,0,2,0)$.
  \item $V$: $(1,1)$, $(0,1)$ y su suma $(1,2)$.
\end{enumerate}
\resultado
\[
Q=\begin{pmatrix}1&1&0&0\\0&0&1&1\\1&1&1&1\end{pmatrix},\qquad
K=\begin{pmatrix}2&0&0&0\\0&0&2&0\\2&0&2&0\end{pmatrix},\qquad
V=\begin{pmatrix}1&1\\0&1\\1&2\end{pmatrix},
\]
1~p cada matriz.
\comprobacion en las tres matrices, la fila~3 es la suma de las filas 1 y~2, porque $x_3=x_1+x_2$.

\apartado{b}{3}
\datos $Q$ y $K$ del apartado~(a); $d_k=4$.
\planteamiento $S_{ij}=q_i\cdot k_j$; después se divide por $\sqrt{d_k}=2$.
\pasos
\begin{enumerate}[label=\arabic*.]
  \item Fila 1, $q_1=(1,1,0,0)$: $q_1\cdot k_1=2$, $q_1\cdot k_2=0$, $q_1\cdot k_3=2$.
  \item Fila 2, $q_2=(0,0,1,1)$: $q_2\cdot k_1=0$, $q_2\cdot k_2=2$, $q_2\cdot k_3=2$.
  \item Fila 3, $q_3=(1,1,1,1)$: $q_3\cdot k_1=2$, $q_3\cdot k_2=2$, $q_3\cdot k_3=4$.
\end{enumerate}
\resultado
\[
S=\begin{pmatrix}2&0&2\\0&2&2\\2&2&4\end{pmatrix},\qquad
\frac{S}{\sqrt{d_k}}=\begin{pmatrix}1&0&1\\0&1&1\\1&1&2\end{pmatrix};
\]
2~p por $S$ y 1~p por el escalado.
\comprobacion la columna~3 de $S$ es la suma de las columnas 1 y~2, porque $k_3=k_1+k_2$.

\apartado{c}{2}
\datos tres posiciones; máscara causal.
\planteamiento $M_{ij}=0$ si $j\le i$ y $M_{ij}=-\infty$ si $j>i$ (M13, §3.1).
\pasos
\begin{enumerate}[label=\arabic*.]
  \item $M$ es triangular superior estricta de $-\infty$; se suma elemento a elemento a $S/\sqrt{d_k}$:
  \[
  M=\begin{pmatrix}0&-\infty&-\infty\\0&0&-\infty\\0&0&0\end{pmatrix},\qquad
  Z=\begin{pmatrix}1&-\infty&-\infty\\0&1&-\infty\\1&1&2\end{pmatrix}.
  \]
  \item Con $-\infty$, $e^{-\infty}=0$: el peso de cada token futuro es exactamente nulo y los pesos permitidos se renormalizan solos. Con~$0$, $e^{0}=1$ y el futuro conservaría peso; por ejemplo, la fila~1 sería $\softmax(1,0,1)=(0{,}422;\allowbreak\ 0{,}155;\allowbreak\ 0{,}422)$.
\end{enumerate}
\resultado \textbf{$M$ y $Z$ como arriba; $-\infty$ anula exactamente los pesos futuros}, mientras que $0$ no anula nada.
\comprobacion la fila~$i$ de $Z$ conserva $i$ valores finitos: 1, 2 y~3.

\apartado{d}{3}
\datos $Z$ del apartado~(c); $e\approx 2{,}718$, $e^{2}\approx 7{,}389$; $V$ del apartado~(a).
\planteamiento softmax por filas con $e^{-\infty}=0$; después $O=AV$.
\pasos
\begin{enumerate}[label=\arabic*.]
  \item Fila 1: solo $z_{11}$ es finito, luego $a_1=(1,0,0)$.
  \item Fila 2: $e^{0}=1$ y $e^{1}=2{,}718$, con suma $3{,}718$: $a_{21}=0{,}269$ y $a_{22}=0{,}731$.
  \item Fila 3: dividiendo $e^{1},e^{1},e^{2}$ por $e$ quedan $1,1,e$, con suma $4{,}718$: $a_{31}=a_{32}=0{,}212$ y $a_{33}=0{,}576$.
  \item $o_1=(1,1)$; $o_2=0{,}269\,(1,1)+0{,}731\,(0,1)=(0{,}269;\allowbreak\ 1)$.
  \item $o_3=0{,}212\,(1,1)+0{,}212\,(0,1)+0{,}576\,(1,2)=(0{,}788;\allowbreak\ 1{,}576)$, con los pesos a cuatro cifras ($0{,}2119$ y $0{,}5761$).
\end{enumerate}
\resultado
\[
A=\begin{pmatrix}1&0&0\\0{,}269&0{,}731&0\\0{,}212&0{,}212&0{,}576\end{pmatrix},\qquad
O=AV=\begin{pmatrix}1&1\\0{,}269&1\\0{,}788&1{,}576\end{pmatrix}.
\]
\comprobacion cada fila de $A$ suma~1; como la segunda componente de $v_1$ y de $v_2$ es~1, la de $o_2$ vale exactamente~1. Criterios: máscara aplicada antes de la softmax (1~p), $A$ correcta (1~p), $O$ correcta (1~p).

\apartado{e}{2}
\datos $x_3=(5,-3)$; el resto, igual.
\planteamiento $A$ es triangular inferior: $o_1$ y $o_2$ solo dependen de los tokens 1 y~2 (M13, §3.2). Solo hay que recalcular la fila~3.
\pasos
\begin{enumerate}[label=\arabic*.]
  \item $o_1$ y $o_2$ no dependen de $x_3$: solo cambia $o_3$.
  \item $q_3=5\,(1,1,0,0)-3\,(0,0,1,1)=(5,5,-3,-3)$; $k_3=5\,(2,0,0,0)-3\,(0,0,2,0)=(10,0,-6,0)$; $v_3=5\,(1,1)-3\,(0,1)=(5,2)$.
  \item Puntuaciones escaladas: $q_3\cdot k_1/2=10/2=5$; $q_3\cdot k_2/2=-6/2=-3$; $q_3\cdot k_3/2=(50+18)/2=34$.
  \item $e^{34}/e^{5}=e^{29}\approx 3{,}9\cdot 10^{12}$, luego $a_{33}\approx 1$, $a_{31}\approx e^{-29}\approx 2{,}5\cdot 10^{-13}$ y $a_{32}\approx e^{-37}\approx 8{,}5\cdot 10^{-17}$.
  \item $o_3\approx v_3=(5,2)$.
\end{enumerate}
\resultado \textbf{solo cambia $o_3$, que pasa a ser $(5{,}000;\allowbreak\ 2{,}000)$}: la cabeza \emph{copia} el valor $v_3$, porque la puntuación de $k_3$ es mucho mayor que las demás.
\comprobacion con Python, $1-a_{33}=2{,}5\cdot 10^{-13}$. Es la situación de «copia en la atención» de la tarea~3 de CS224N~\cite{cs224n}: casi todo el peso recae en un solo elemento cuando $q\cdot k_j\gg q\cdot k_i$ para todo $i\ne j$.

\apartado{f}{2}
\datos $a_{22}=0{,}731$, del apartado~(d).
\planteamiento para la softmax, $\partial a_{ij}/\partial z_{ij}=a_{ij}(1-a_{ij})$; el escalado por $\sqrt{d_k}$ evita la saturación~\cite[§3.2.1]{vaswani} (M13, §2.3).
\pasos
\begin{enumerate}[label=\arabic*.]
  \item $\partial a_{22}/\partial z_{22}=0{,}731\cdot 0{,}269=0{,}197$.
  \item Si las puntuaciones crecen sin escalar, la softmax se satura: $a_{ij}\to 0$ o $a_{ij}\to 1$, y en ambos casos $a_{ij}(1-a_{ij})\to 0$.
  \item Sin gradiente, la cabeza deja de aprender a quién atender.
\end{enumerate}
\resultado \textbf{$\partial a_{22}/\partial z_{22}=0{,}197$; con puntuaciones grandes tiende a~0 (softmax saturada).}
\comprobacion el máximo de $a(1-a)$ es $0{,}25$, en $a=0{,}5$, y $0{,}197<0{,}25$. Respecto de la puntuación sin escalar $s=2z$, la regla de la cadena añade el factor $1/\sqrt{d_k}=1/2$: $0{,}098$.

\repasar{M13, §2.1, §2.3, §3.1 y §3.2}
\end{solucion}

% ===================================================================== B3
\begin{solucion}{eje:B-bpe}

\apartado{a}{1}
\datos \tok{pino}~$\times 4$, \tok{pinos}~$\times 2$, \tok{vino}~$\times 3$, \tok{vinos}~$\times 1$, \tok{pinta}~$\times 2$.
\planteamiento número de símbolos $=\sum(\text{longitud}\times\text{frecuencia})$.
\pasos
\begin{enumerate}[label=\arabic*.]
  \item $4\cdot 4+5\cdot 2+4\cdot 3+5\cdot 1+5\cdot 2=16+10+12+5+10$.
\end{enumerate}
\resultado \textbf{53 símbolos.}
\comprobacion hay $4+2+3+1+2=12$ palabras de 4 o 5 letras: el total debe estar entre 48 y~60.

\apartado{b}{6}
\datos el corpus del apartado~(a), partido en caracteres.
\planteamiento en cada paso se cuentan los pares adyacentes ponderados por la frecuencia de la palabra, se fusiona el más frecuente y se restan sus apariciones del total de símbolos~\cite{sennrich} (M12, §2.1). 1,5~p por paso.
\pasos
\begin{enumerate}[label=\arabic*.]
  \item Recuento inicial: $(\tok{i},\tok{n})=4+2+3+1+2=12$; $(\tok{n},\tok{o})=4+2+3+1=10$; $(\tok{p},\tok{i})=4+2+2=8$; $(\tok{v},\tok{i})=4$; $(\tok{o},\tok{s})=3$. Gana $(\tok{i},\tok{n})$; símbolos: $53-12=41$.
  \item Tras \tok{in}: $(\tok{in},\tok{o})=10$; $(\tok{p},\tok{in})=8$; $(\tok{v},\tok{in})=4$; $(\tok{o},\tok{s})=3$. Gana $(\tok{in},\tok{o})$; símbolos: $41-10=31$.
  \item Tras \tok{ino}: $(\tok{p},\tok{ino})=4+2=6$; $(\tok{v},\tok{ino})=4$; $(\tok{ino},\tok{s})=3$. Gana $(\tok{p},\tok{ino})$; símbolos: $31-6=25$.
  \item Tras \tok{pino}: $(\tok{v},\tok{ino})=4$; empatan con~2 $(\tok{pino},\tok{s})$, $(\tok{p},\tok{in})$, $(\tok{in},\tok{t})$ y $(\tok{t},\tok{a})$. Gana $(\tok{v},\tok{ino})$; símbolos: $25-4=21$.
\end{enumerate}

\begin{center}\small
\begin{tabularx}{\textwidth}{@{}cll>{\raggedright\arraybackslash}Xc@{}}
\toprule
Paso & Par (frec.) & Segundo (frec.) & Corpus tras la fusión & Símbolos \\
\midrule
1 & \tok{i n} (12) & \tok{n o} (10) & \tok{p in o}$\times4$, \tok{p in o s}$\times2$, \tok{v in o}$\times3$, \tok{v in o s}$\times1$, \tok{p in t a}$\times2$ & 41 \\
2 & \tok{in o} (10) & \tok{p in} (8) & \tok{p ino}$\times4$, \tok{p ino s}$\times2$, \tok{v ino}$\times3$, \tok{v ino s}$\times1$, \tok{p in t a}$\times2$ & 31 \\
3 & \tok{p ino} (6) & \tok{v ino} (4) & \tok{pino}$\times4$, \tok{pino s}$\times2$, \tok{v ino}$\times3$, \tok{v ino s}$\times1$, \tok{p in t a}$\times2$ & 25 \\
4 & \tok{v ino} (4) & cuatro pares (2) & \tok{pino}$\times4$, \tok{pino s}$\times2$, \tok{vino}$\times3$, \tok{vino s}$\times1$, \tok{p in t a}$\times2$ & 21 \\
\bottomrule
\end{tabularx}
\end{center}

\resultado \textbf{reglas, en orden: \tok{in}, \tok{ino}, \tok{pino}, \tok{vino}; el corpus queda con 21 símbolos.}
\comprobacion recuento directo: $1\cdot 4+2\cdot 2+1\cdot 3+2\cdot 1+4\cdot 2=21$. No hay empates en el primer puesto, de modo que la regla de desempate (M12, §2.3) no interviene. Error frecuente: tomar $(\tok{p},\tok{i})=8$ como máximo olvidando las apariciones en \tok{vino}.

\apartado{c}{3}
\datos reglas en orden: $(\tok{i},\tok{n})$, $(\tok{in},\tok{o})$, $(\tok{p},\tok{ino})$, $(\tok{v},\tok{ino})$.
\planteamiento partir cada palabra en caracteres y aplicar las reglas en el orden aprendido (M12, §2.2).
\pasos
\begin{enumerate}[label=\arabic*.]
  \item \tok{vinos}: \tok{v i n o s} $\to$ \tok{v in o s} $\to$ \tok{v ino s} $\to$ \tok{vino s}: 2 tokens.
  \item \tok{pinto}: \tok{p i n t o} $\to$ \tok{p in t o}; la regla~2 no se aplica (tras \tok{in} viene \tok{t}) y la~3 necesita \tok{ino}: 4 tokens.
  \item \tok{sino}: \tok{s i n o} $\to$ \tok{s in o} $\to$ \tok{s ino}; no existe la regla $(\tok{s},\tok{ino})$: 2 tokens.
  \item \tok{ovino}: \tok{o v i n o} $\to$ \tok{o v in o} $\to$ \tok{o v ino} $\to$ \tok{o vino}: 2 tokens.
\end{enumerate}
\resultado \textbf{\tok{vinos} $=$ [\tok{vino}, \tok{s}] (2); \tok{pinto} $=$ [\tok{p}, \tok{in}, \tok{t}, \tok{o}] (4); \tok{sino} $=$ [\tok{s}, \tok{ino}] (2); \tok{ovino} $=$ [\tok{o}, \tok{vino}] (2).}
\comprobacion ninguna palabra nueva produce un token desconocido, porque los caracteres sueltos pertenecen siempre al vocabulario.

\apartado{d}{2}
\datos «piñón» en UTF-8.
\planteamiento los caracteres ASCII ocupan 1~byte; «ñ» y «ó», 2 (M12, §1.3 y §1.5).
\pasos
\begin{enumerate}[label=\arabic*.]
  \item Bytes en hexadecimal: \texttt{70}, \texttt{69}, \texttt{C3\,B1}, \texttt{C3\,B3}, \texttt{6E}: $1+1+2+2+1=7$.
  \item En español, los pares $(\texttt{C3},\texttt{B1})$, que forman «ñ», y $(\texttt{C3},\texttt{B3})$, que forma «ó» (y análogos como $(\texttt{C3},\texttt{A1})$ para «á»), son muy frecuentes.
  \item Sin esas fusiones, ninguna palabra con tilde o con «ñ» puede llegar a ser un único token.
\end{enumerate}
\resultado \textbf{7 símbolos iniciales; conviene fusionar pronto $(\tok{C3},\tok{B1})$ y $(\tok{C3},\tok{B3})$}, los dos bytes de «ñ» y de «ó».
\comprobacion coincide con la cuestión~\ref{q:pinon} (7~bytes).

\repasar{M12, §1.3, §1.5, §2.1, §2.2 y §2.4}
\end{solucion}

% ===================================================================== B4
\begin{solucion}{eje:B-lince}

\apartado{a}{2}
\datos $d=2048$.
\planteamiento $\dff=\tfrac{8}{3}d$ redondeado hacia arriba al múltiplo de 256~\cite{shazeer} (M14, §3.3).
\pasos
\begin{enumerate}[label=\arabic*.]
  \item $8\cdot 2048/3=16\,384/3=5461{,}33$.
  \item $5461{,}33/256=21{,}33$; el entero siguiente es~22.
  \item $22\cdot 256=5632$.
\end{enumerate}
\resultado \textbf{$\dff=5632$.}
\comprobacion $21\cdot 256=5376<5461{,}33<5632$.

\apartado{b}{4}
\datos $d=2048$, $L=24$, $\dff=5632$, $V=48\,000$; embeddings atados; RMSNorm solo con pesos; sin sesgos.
\planteamiento $N=L\,(4d^{2}+3d\,\dff+2d)+Vd+d$, con una sola matriz de embeddings por estar atados (M14, §5.1 y §5.2).
\pasos
\begin{enumerate}[label=\arabic*.]
  \item $d^{2}=4\,194\,304$; atención: $4d^{2}=16\,777\,216$.
  \item MLP: $3\cdot 2048\cdot 5632=34\,603\,008$.
  \item Dos RMSNorm: $2\cdot 2048=4096$. Bloque: $16\,777\,216+34\,603\,008+4096=51\,384\,320$.
  \item 24 bloques: $24\cdot 51\,384\,320=1\,233\,223\,680$.
  \item Embedding: $48\,000\cdot 2048=98\,304\,000$; RMSNorm final: $2048$.
  \item $N=1\,233\,223\,680+98\,304\,000+2048=1\,331\,529\,728$.
  \item $12Ld^{2}=12\cdot 24\cdot 4\,194\,304=1\,207\,959\,552$, que es el $90{,}7\,\%$ de $N$.
  \item MLP: $24\cdot 34\,603\,008=830\,472\,192$, el $62{,}4\,\%$ de $N$.
\end{enumerate}

\begin{center}\small
\begin{tabular}{@{}llr@{}}
\toprule
Pieza & Fórmula & Parámetros \\
\midrule
Atención & $4d^{2}$ & $16\,777\,216$ \\
MLP SwiGLU & $3d\,\dff$ & $34\,603\,008$ \\
Dos RMSNorm & $2d$ & $4096$ \\
\midrule
\textbf{Bloque} & & $\mathbf{51\,384\,320}$ \\
\bottomrule
\end{tabular}
\end{center}

\resultado \textbf{bloque: $51\,384\,320$; $N=1\,331\,529\,728\approx 1{,}33\cdot 10^{9}$; $12Ld^{2}$ es el $90{,}7\,\%$ de $N$; los MLP suponen el $62{,}4\,\%$.}
\comprobacion la diferencia con la aproximación procede del embedding ($98\,304\,000/N=7{,}4\,\%$) y del redondeo de $\dff$ hacia arriba. Sin atar, $N$ sería $1\,429\,833\,728$.

\apartado{c}{3}
\datos $N=1{,}3315\cdot 10^{9}$; $D=2\cdot 10^{11}$; H100 a $989\cdot 10^{12}$~FLOP/s; aprovechamiento del 40\,\%; 64~GPU.
\planteamiento inferencia, $\approx 2N$ FLOPs por token; entrenamiento, $C\approx 6ND$~\cite{kaplan}; tiempo $=C/\text{velocidad efectiva}$ (M14, §6.2 y §6.3).
\pasos
\begin{enumerate}[label=\arabic*.]
  \item Inferencia: $2\cdot 1{,}3315\cdot 10^{9}=2{,}663\cdot 10^{9}$ FLOPs por token.
  \item Entrenamiento: $C=6\cdot 1{,}3315\cdot 10^{9}\cdot 2\cdot 10^{11}=1{,}598\cdot 10^{21}$ FLOPs.
  \item Velocidad efectiva por GPU: $0{,}4\cdot 989\cdot 10^{12}=3{,}956\cdot 10^{14}$ FLOP/s.
  \item Tiempo: $1{,}598\cdot 10^{21}/3{,}956\cdot 10^{14}=4{,}039\cdot 10^{6}$~s, es decir, $1122$~h de GPU.
  \item Con 64~GPU: $1122/64=17{,}53$~h de reloj ($0{,}73$~días).
\end{enumerate}
\resultado \textbf{$2{,}66\cdot 10^{9}$ FLOPs por token; $C\approx 1{,}60\cdot 10^{21}$ FLOPs; unas $1122$~h de H100; $17{,}5$~h de reloj con 64~GPU.}
\comprobacion LLaMA-2~7B~\cite{llama2} ($N=6{,}74\cdot 10^{9}$, $D=2\cdot 10^{12}$) cuesta $8{,}09\cdot 10^{22}$ FLOPs, unas $50{,}6$ veces más, coherente con un $N$ cinco veces mayor y un $D$ diez veces mayor.

\apartado{d}{3}
\datos $N=1\,331\,529\,728$; bf16, 2~bytes; int4, $0{,}5$~bytes; Adam en precisión mixta, 16~bytes por parámetro.
\planteamiento memoria $=N\times$ bytes por parámetro (M14, §5.3).
\pasos
\begin{enumerate}[label=\arabic*.]
  \item bf16: $2N=2{,}663\cdot 10^{9}$~bytes $=2{,}66$~GB.
  \item int4: $0{,}5N=6{,}658\cdot 10^{8}$~bytes $=0{,}67$~GB.
  \item Adam en precisión mixta: $16N=2{,}130\cdot 10^{10}$~bytes $=21{,}3$~GB.
  \item $21{,}3<24$: los estados caben, pero solo quedan unos $2{,}7$~GB para las activaciones.
\end{enumerate}
\resultado \textbf{$2{,}66$~GB en bf16, $0{,}67$~GB en int4 y $21{,}3$~GB para el estado del entrenamiento.} El estado cabe en 24~GB, pero las activaciones, que crecen con el lote y el contexto, no caben con holgura.
\comprobacion $21{,}3/2{,}66=8=16/2$. En la práctica se recurriría a la recomputación de activaciones o al reparto entre varias GPU (ZeRO, módulo~17).

\apartado{e}{3}
\datos $L=24$, $d=2048$, $n=4096$; $2N=2{,}663\cdot 10^{9}$.
\planteamiento comparar $4Lnd$ con $2N$ y despejar $n$ en la igualdad: $n^{*}=N/(2Ld)$ (M14, §6.2).
\pasos
\begin{enumerate}[label=\arabic*.]
  \item $4\cdot 24\cdot 4096\cdot 2048=805\,306\,368\approx 8{,}053\cdot 10^{8}$ FLOPs por token.
  \item Fracción: $8{,}053\cdot 10^{8}/2{,}663\cdot 10^{9}=0{,}302$.
  \item Igualdad: $n^{*}=2N/(4Ld)=2{,}663\cdot 10^{9}/196\,608=13\,545$.
\end{enumerate}
\resultado \textbf{la atención añade un $30{,}2\,\%$ a $2N$ con $n=4096$, y lo iguala con $n\approx 13\,545$ tokens.}
\comprobacion la fracción es lineal en $n$: $4096/13\,545=0{,}302$.

\repasar{M14, §3.3, §5.1, §5.2, §5.3, §6.2 y §6.3}
\end{solucion}

% ===================================================================== B5
\begin{solucion}{eje:B-rope}

\apartado{a}{5}
\datos $d_k=2$, $\theta=\pi/4$, $q=(1,0)$, $k=(1,1)$; $\cos(\pi/4)=\sin(\pi/4)=0{,}707$.
\planteamiento calcular las dos constantes de la identidad y evaluar en cada $\varphi=(n-m)\theta$~\cite{roformer} (M14, §4.3).
\pasos
\begin{enumerate}[label=\arabic*.]
  \item $q\cdot k=1$ y $q_1k_2-q_2k_1=1$; la puntuación es $\cos\varphi-\sin\varphi$.
  \item $(2,1)$ y $(5,4)$: $n-m=-1$, $\varphi=-\pi/4$, puntuación $0{,}707+0{,}707=1{,}414$.
  \item $(3,3)$: $\varphi=0$, puntuación $1$.
  \item $(1,2)$: $\varphi=\pi/4$, puntuación $0$.
  \item $(0,4)$: $\varphi=\pi$, puntuación $-1$.
\end{enumerate}

\begin{center}\small
\begin{tabular}{@{}cccc@{}}
\toprule
$(m,n)$ & $n-m$ & $\varphi$ & $\cos\varphi-\sin\varphi$ \\
\midrule
$(2,1)$ & $-1$ & $-\pi/4$ & $1{,}414$ \\
$(5,4)$ & $-1$ & $-\pi/4$ & $1{,}414$ \\
$(3,3)$ & $0$ & $0$ & $1$ \\
$(1,2)$ & $+1$ & $\pi/4$ & $0$ \\
$(0,4)$ & $+4$ & $\pi$ & $-1$ \\
\bottomrule
\end{tabular}
\end{center}

\resultado \textbf{$1{,}414$; $1{,}414$; $1$; $0$; $-1$.} Los dos primeros pares ocupan posiciones distintas a la misma distancia y dan la misma puntuación: RoPE codifica la posición relativa.
\comprobacion con distancia nula se recupera $q\cdot k=1$; además, $\cos\varphi-\sin\varphi=\sqrt2\cos(\varphi+\pi/4)$, cuyo máximo, $\sqrt2=1{,}414$, se alcanza en $\varphi=-\pi/4$. Criterios: 1~p por las dos primeras filas, 1~p por cada una de las otras tres; la observación es necesaria para obtener la puntuación completa.

\apartado{b}{2}
\datos $d=2048$, $h=16$.
\planteamiento $d_k=d/h$; los parámetros de la atención son $4d^{2}$, sin dependencia de $h$ (M13, §4.2).
\pasos
\begin{enumerate}[label=\arabic*.]
  \item $d_k=2048/16=128$.
  \item Con $h=32$, $d_k=64$, pero las 32 cabezas siguen formando $W_Q$, $W_K$, $W_V$ y $W_O$ de tamaño $d\times d$: $4d^{2}=16\,777\,216$ en ambos casos.
\end{enumerate}
\resultado \textbf{$d_k=128$; duplicar $h$ no añade parámetros.}
\comprobacion coincide con la cuestión~\ref{q:mha}: el recuento $4d^{2}$ no depende de~$h$.

\apartado{c}{4}
\datos $n=32\,768$; fp16, 2~bytes; $h=16$; $L=24$; $d=2048$.
\planteamiento cada matriz $A$ es $n\times n$; FLOPs cuadráticos $4n^{2}d$ frente a lineales $8nd^{2}$ (M13, §4.3 y §4.4).
\pasos
\begin{enumerate}[label=\arabic*.]
  \item Una cabeza: $n^{2}=1\,073\,741\,824$ elementos; por 2~bytes, $2{,}147\cdot 10^{9}$~bytes $=2{,}15$~GB.
  \item Todas las cabezas y capas: $2{,}147\cdot 10^{9}\cdot 16\cdot 24=8{,}246\cdot 10^{11}$~bytes $\approx 825$~GB.
  \item Cruce: $4n^{2}d>8nd^{2}$ equivale, dividiendo por $4nd>0$, a $n>2d=4096$.
  \item FlashAttention calcula la softmax por bloques y nunca almacena $A$ completa: la memoria adicional es lineal en $n$, aunque el número de FLOPs sigue siendo $O(n^{2}d)$~\cite{dao}.
\end{enumerate}
\resultado \textbf{$2{,}15$~GB por cabeza y unos 825~GB en total; la atención domina en FLOPs a partir de $n=4096$; FlashAttention evita almacenar~$A$.}
\comprobacion 825~GB supera con mucho los 80~GB de una GPU actual: sin FlashAttention, ese contexto es inviable.

\repasar{M13, §4.2, §4.3 y §4.4; M14, §4.3}
\end{solucion}

% ===================================================================== B6
\begin{solucion}{eje:B-minigpt}

\apartado{a}{3}
\datos $V=80$, $T=128$, $d=128$, $L=4$; MLP de $4d$; LayerNorm con peso y sesgo; posiciones aprendidas; embeddings atados; lineales sin sesgo.
\planteamiento bloque: $4d^{2}+8d^{2}+2\cdot 2d$; fuera de los bloques: $Vd$ (una vez), $Td$ y la LayerNorm final, $2d$ (M15, §2.3).
\pasos
\begin{enumerate}[label=\arabic*.]
  \item $d^{2}=16\,384$; atención: $65\,536$; MLP: $8d^{2}=131\,072$.
  \item Dos LayerNorm: $2\cdot 2\cdot 128=512$. Bloque: $197\,120$.
  \item Cuatro bloques: $788\,480$.
  \item Embedding: $80\cdot 128=10\,240$; posiciones: $128\cdot 128=16\,384$; LayerNorm final: $256$.
  \item $N=788\,480+10\,240+16\,384+256=815\,360$.
\end{enumerate}
\resultado \textbf{$N=815\,360$ parámetros.}
\comprobacion $12Ld^{2}=786\,432$, el $96{,}5\,\%$ de $N$.

\apartado{b}{2}
\datos $V=80$; pérdida inicial $\approx 4{,}39$.
\planteamiento logits casi iguales dan una softmax casi uniforme y una pérdida $\ln V$ (M15, §1.4).
\pasos
\begin{enumerate}[label=\arabic*.]
  \item Con pesos pequeños, los logits son casi iguales y $p_i\approx 1/80$.
  \item $-\ln(1/80)=\ln 80=\ln 8+\ln 10=2{,}079+2{,}303=4{,}382$.
\end{enumerate}
\resultado \textbf{$\ln 80=4{,}382\approx 4{,}39$.}
\comprobacion una pérdida inicial mucho mayor indicaría una inicialización defectuosa, por ejemplo embeddings $\mathcal N(0,1)$ atados.

\apartado{c}{3}
\datos \cref{tab:curvas}; bigramas, $2{,}45$; $\ln 2=0{,}693$.
\planteamiento se elige el mínimo de validación; $\PPL=e^{\ell}$ y bits por carácter $=\ell/\ln 2$ (M15, §1.3 y §4.2).
\pasos
\begin{enumerate}[label=\arabic*.]
  \item Mínimo de validación: paso 2500, con $1{,}44$ (antes, $1{,}46$; después, $1{,}45$).
  \item $\PPL=e^{1{,}44}=4{,}221$.
  \item Bits por carácter: $1{,}44/0{,}693=2{,}078$.
  \item Bigramas: $e^{2{,}45}=11{,}59$; cociente $11{,}59/4{,}221=2{,}75$.
\end{enumerate}
\resultado \textbf{punto de control del paso 2500; $\PPL\approx 4{,}22$ y $2{,}08$ bits por carácter; el modelo de bigramas tiene $\PPL\approx 11{,}59$}: el mini-GPT duda entre casi tres veces menos opciones.
\comprobacion $2^{2{,}078}=4{,}22$: la perplejidad no depende de la base del logaritmo.

\apartado{d}{2}
\datos $B=32$, $T=128$, 5000~pasos; corpus de $10^{6}$ caracteres.
\planteamiento tokens $=B\cdot T\cdot\text{pasos}$; épocas $=$ tokens $/$ tamaño del corpus (M15, §1.2 y §4.4).
\pasos
\begin{enumerate}[label=\arabic*.]
  \item $32\cdot 128=4096$ tokens por paso; por 5000 pasos, $20\,480\,000$ tokens.
  \item Épocas: $20\,480\,000/10^{6}=20{,}48$.
  \item Brecha entre validación y entrenamiento: $1{,}44-1{,}30=0{,}14$ en el paso 2500 y $1{,}62-0{,}98=0{,}64$ en el 5000.
  \item El modelo ve cada carácter unas 20 veces y memoriza: la validación sube desde el paso 2500.
\end{enumerate}
\resultado \textbf{$20\,480\,000$ tokens, unas $20{,}5$ épocas: sobreajuste.} Remedios: más datos (el mejor), parada temprana, \emph{dropout}, más decaimiento de pesos o un modelo más pequeño.
\comprobacion la pérdida de entrenamiento baja mientras la de validación sube: es la firma del sobreajuste en la tabla de síntomas de M15, §4.4.

\apartado{e}{2}
\datos logits $(2,1,0)$; $\tau=1$ y $\tau=0{,}5$.
\planteamiento $p=\softmax(z/\tau)$ (M15, §5.2).
\pasos
\begin{enumerate}[label=\arabic*.]
  \item $\tau=1$: $e^{2}=7{,}389$, $e^{1}=2{,}718$, $e^{0}=1$, con suma $11{,}107$: $p=(0{,}665;\allowbreak\ 0{,}245;\allowbreak\ 0{,}090)$.
  \item $\tau=0{,}5$: $z/\tau=(4,2,0)$; $e^{4}=54{,}60$, $e^{2}=7{,}389$, $e^{0}=1$, con suma $62{,}99$: $p=(0{,}867;\allowbreak\ 0{,}117;\allowbreak\ 0{,}016)$.
  \item Bajar $\tau$ concentra la masa en el carácter más probable: las muestras son más correctas, pero más repetitivas. Subir $\tau$ (con $\tau=2$, $p=(0{,}506;\allowbreak\ 0{,}307;\allowbreak\ 0{,}186)$) da más variedad y más errores.
\end{enumerate}
\resultado \textbf{$\tau=1$: $(0{,}665;\allowbreak\ 0{,}245;\allowbreak\ 0{,}090)$; $\tau=0{,}5$: $(0{,}867;\allowbreak\ 0{,}117;\allowbreak\ 0{,}016)$.} Bajar la temperatura hace las muestras más predecibles y repetitivas.
\comprobacion ambas distribuciones suman~1 y, cuando $\tau\to 0^{+}$, la distribución tiende a $(1,0,0)$.

\repasar{M15, §1.2, §1.3, §1.4, §2.3, §4.2, §4.4 y §5.2}
\end{solucion}

% =====================================================================
\chapter{Autoevaluación}
\label{ap:autoevaluacion}
% =====================================================================

\section{Cálculo de la nota}

La nota es la suma de los puntos de la parte~A y de la parte~B, sobre 100. El examen se aprueba con \textbf{70 puntos} (70\,\%). Aunque se apruebe, si en alguna parte o en algún módulo se obtiene menos del 50\,\% de los puntos posibles, debe repasarse ese contenido antes de continuar.

\begin{table}[ht]
\centering\small
\caption{Puntuación por parte.}\label{tab:partes}
\begin{tabular}{@{}lllcc@{}}
\toprule
Parte & Contenido & Módulos & Puntos & Obtenidos \\
\midrule
A & Cuestiones (20 $\times$ 1~p) & M12--M15 & 20 & \\
B.1 & Bloque decoder & M14 & 15 & \\
B.2 & Atención con máscara causal & M13 & 15 & \\
B.3 & BPE a mano & M12 & 12 & \\
B.4 & Parámetros, FLOPs y memoria & M14 & 15 & \\
B.5 & RoPE y coste de la atención & M13, M14 & 11 & \\
B.6 & Entrenamiento del mini-GPT & M15 & 12 & \\
\midrule
\textbf{Total} & & & \textbf{100} & \\
\bottomrule
\end{tabular}
\end{table}

\begin{table}[ht]
\centering\small
\caption{Puntuación por módulo. En el ejercicio B.5, el apartado~(a) corresponde a M14 y los apartados (b) y~(c), a M13.}\label{tab:modulos}
\begin{tabular}{@{}llccc@{}}
\toprule
Módulo & Elementos & Puntos & 50\,\% & Obtenidos \\
\midrule
M12. Tokenización & A.1--A.5; B.3 & $5+12=17$ & $8{,}5$ & \\
M13. Atención & A.6--A.10; B.2; B.5(b,c) & $5+15+6=26$ & $13$ & \\
M14. Transformer & A.11--A.18; B.1; B.4; B.5(a) & $8+15+15+5=43$ & $21{,}5$ & \\
M15. Mini-GPT & A.19--A.20; B.6 & $2+12=14$ & $7$ & \\
\midrule
\textbf{Total} & & \textbf{100} & & \\
\bottomrule
\end{tabular}
\end{table}

\section{Puerta de salida de la fase III}

Para pasar a la fase~IV deben cumplirse las tres condiciones del método de estudio del curso:
\begin{enumerate}[label=\arabic*.]
  \item en el test de cada módulo (M12, M13, M14 y M15), al menos 9 aciertos de 12;
  \item en la autoevaluación del laboratorio interactivo de la fase, al menos un 70\,\%;
  \item en este examen de hito, al menos 70 puntos de 100.
\end{enumerate}
El proyecto de fase se considera superado con al menos 70 puntos en la rúbrica del \cref{tab:rubrica}.

\section{Qué hacer si no se alcanza el umbral}

\begin{itemize}[leftmargin=1.4em,label=\textendash]
  \item Localizar en el \cref{tab:repaso} las secciones correspondientes a cada fallo y anotar el error en el registro de errores, clasificado como de concepto, de cálculo o de despiste.
  \item Releer esas secciones y rehacer, sin consultar la solución, los ejercicios $\star\star$ y $\star\star\star$ del tema correspondiente de los apuntes del módulo.
  \item Repetir el test del módulo afectado hasta alcanzar 9 de 12 y volver a resolver, al cabo de unos días, los ejercicios de la parte~B en los que se perdió más de la mitad de los puntos.
  \item Si la nota está por debajo de 50, repasar la fase completa siguiendo el ciclo del método de estudio (teoría, laboratorio, práctica, ejercicios, test y tarjetas) antes de repetir el examen.
\end{itemize}

\begin{table}[ht]
\centering\small
\caption{Secciones que deben repasarse según los fallos.}\label{tab:repaso}
\begin{tabularx}{\textwidth}{@{}>{\raggedright\arraybackslash}p{4.3cm}X@{}}
\toprule
Fallos & Repasar \\
\midrule
A.1, A.4; B.3(b,c) & M12, §2.1 (entrenamiento de BPE) y §2.2 (codificación) \\
A.2, A.3; B.3(d) & M12, §1.3 (UTF-8), §1.5 (consecuencias para el español) y §2.4 (BPE a nivel de byte) \\
A.5 & M12, §5.5 (secuencialidad y paso al Transformer) \\
A.6; B.2(a,b,d) & M13, §2.1 (forma matricial de la atención) \\
A.7; B.2(f) & M13, §2.3 (por qué se divide por $\sqrt{d_k}$) \\
A.8; B.2(c,e) & M13, §3.1 (máscaras aditivas) y §3.2 (causalidad) \\
A.9; B.5(b) & M13, §4.2 (parámetros de la atención multi-cabeza) \\
A.10; B.5(c) & M13, §4.3 (coste en operaciones) y §4.4 (memoria y contexto largo) \\
B.1(a--d) & M14, §1.2 (bloque decoder-only), §2.1 (flujo residual) y §3.1 (MLP posición a posición) \\
A.11, A.12; B.1(e) & M14, §2.2 (pre-norm frente a post-norm) y §2.3 (crecimiento del flujo) \\
A.13; B.4(a) & M14, §3.3 (SwiGLU) \\
A.14, A.15; B.5(a) & M14, §4.3 (RoPE) \\
A.16 & M14, §1.4 (embeddings atados) y §5.2 (parámetros del modelo completo) \\
A.17, A.18; B.4(b--e) & M14, §5.1 a §5.3 (parámetros y memoria) y §6.2 y §6.3 (FLOPs y regla $6ND$) \\
A.19; B.6(a,b) & M15, §1.4 (pérdida de la predicción uniforme) y §2.3 (recuento de parámetros) \\
A.20; B.6(c,d) & M15, §1.3 (perplejidad), §4.2 (puntos de control) y §4.4 (lectura de curvas) \\
B.6(e) & M15, §5.2 (temperatura) \\
\bottomrule
\end{tabularx}
\end{table}

\paragraph{Siguiente paso.} Superada la puerta de salida, se continúa con la fase~IV, módulo~16: preentrenamiento a escala.


\backmatter
% =====================================================================
\begin{thebibliography}{99}
\addcontentsline{toc}{chapter}{Bibliografía}
\markboth{Bibliografía}{}
\small

\bibitem{cs224n}
D.~Yang y Y.~Choi.
\newblock CS224N: Natural Language Processing with Deep Learning.
\newblock Stanford University, invierno de 2026. Tarea~3, «Self-Attention and Transformers»: pregunta~1, «Attention exploration», y pregunta~2, «Position embeddings exploration».
\newblock \url{https://web.stanford.edu/class/cs224n/}

\bibitem{cs336}
Equipo docente de CS336 (T.~Hashimoto y P.~Liang).
\newblock CS336: Language Modeling from Scratch. Tarea~1, «Basics: Building a Transformer LM», versión 26.0.3.
\newblock Stanford University, primavera de 2026. Se citan el §2.4 (ejemplo \emph{bpe\_example}), el §2.6.1 (ejemplo \emph{bpe\_encoding}), los §3.4 y §3.5 (problemas \emph{rope}, \emph{transformer\_block}, \emph{transformer\_lm} y \emph{transformer\_accounting}), el §4.1 (perplejidad), el §4.3 (problema \emph{adamw\_accounting}), el §6 (problema \emph{decoding}) y el §7.3 (problema \emph{pre\_norm\_ablation}).
\newblock \url{https://github.com/stanford-cs336/assignment1-basics}

\bibitem{dao}
T.~Dao, D.~Y. Fu, S.~Ermon, A.~Rudra y C.~Ré.
\newblock FlashAttention: Fast and memory-efficient exact attention with IO-awareness.
\newblock En \emph{Advances in Neural Information Processing Systems 35}, 2022.
\newblock \url{https://arxiv.org/abs/2205.14135}

\bibitem{jm3}
D.~Jurafsky y J.~H. Martin.
\newblock \emph{Speech and Language Processing}, 3.ª ed., borrador del 19 de agosto de 2026.
\newblock Cap.~2, «Words and Tokens» (§2.4.1 y ejercicio~2.3); cap.~7, «Transformers and Pretraining» (ejercicios~7.1, 7.2, 7.4 y~7.5).
\newblock \url{https://web.stanford.edu/~jurafsky/slp3/}

\bibitem{kaplan}
J.~Kaplan, S.~McCandlish, T.~Henighan, T.~B. Brown, B.~Chess, R.~Child, S.~Gray, A.~Radford, J.~Wu y D.~Amodei.
\newblock Scaling laws for neural language models.
\newblock Prepublicación, 2020.
\newblock \url{https://arxiv.org/abs/2001.08361}

\bibitem{gpt2}
A.~Radford, J.~Wu, R.~Child, D.~Luan, D.~Amodei e I.~Sutskever.
\newblock Language models are unsupervised multitask learners.
\newblock Informe técnico, OpenAI, 2019.

\bibitem{sennrich}
R.~Sennrich, B.~Haddow y A.~Birch.
\newblock Neural machine translation of rare words with subword units.
\newblock En \emph{Proceedings of the 54th Annual Meeting of the Association for Computational Linguistics (ACL)}, págs.~1715--1725, 2016.
\newblock \url{https://aclanthology.org/P16-1162/}

\bibitem{shazeer}
N.~Shazeer.
\newblock GLU variants improve Transformer.
\newblock Prepublicación, 2020.
\newblock \url{https://arxiv.org/abs/2002.05202}

\bibitem{roformer}
J.~Su, Y.~Lu, S.~Pan, A.~Murtadha, B.~Wen e Y.~Liu.
\newblock RoFormer: Enhanced Transformer with rotary position embedding.
\newblock \emph{Neurocomputing}, 568, 2024. Prepublicación de 2021.
\newblock \url{https://arxiv.org/abs/2104.09864}

\bibitem{llama2}
H.~Touvron \emph{et al.}
\newblock Llama 2: Open foundation and fine-tuned chat models.
\newblock Prepublicación, 2023.
\newblock \url{https://arxiv.org/abs/2307.09288}

\bibitem{vaswani}
A.~Vaswani, N.~Shazeer, N.~Parmar, J.~Uszkoreit, L.~Jones, A.~N. Gomez, Ł.~Kaiser e I.~Polosukhin.
\newblock Attention is all you need.
\newblock En \emph{Advances in Neural Information Processing Systems 30}, 2017.
\newblock \url{https://arxiv.org/abs/1706.03762}

\bibitem{xiong}
R.~Xiong, Y.~Yang, D.~He, K.~Zheng, S.~Zheng, C.~Xing, H.~Zhang, Y.~Lan, L.~Wang y T.-Y. Liu.
\newblock On layer normalization in the Transformer architecture.
\newblock En \emph{International Conference on Machine Learning (ICML)}, 2020.
\newblock \url{https://arxiv.org/abs/2002.04745}

\end{thebibliography}


\end{document}
